\section{Representatives, compactness, and regularization}
\label{sec:analysis}

We fix one of the finite topological problems constructed in
Section~\ref{sec:stack}, with the metrics, artificial lengths, backgrounds,
and tolerances chosen in Sections~\ref{sec:geometry} and
\ref{sec:estimates}.  Only the actual $J=S^3$ and lens cuts may now have
infinite length.  In particular, a very long $Y$ seam is still a finite
part of this problem.  Constants below may depend on this finite problem.
Connections have the specified integral determinant, transitions are fixed,
and all gauges have determinant one.  The letters $A,R,S,P$ in this section
denote the field types of Definition~\ref{def:indices:types}, not the capped
Seifert surface $A$.

For the spinor regularization, the ordinary equation and representative
data are based at the choices
of Proposition~\ref{prop:stack:composition}, with the zero-winding and
segmentation prescriptions already installed.  Work inside its open
neighborhood of genuine-face-preserving data on which the ordinary
count is $\Omega$.  The weighted perturbation norm below maps
continuously into the finitely many norms defining this neighborhood.
Every subsequent generic choice is made in a sufficiently small relative
open ball satisfying the harmonic-reducer, bit-identification, and
parameter-locality conditions.  These conditions leave the local
variation spaces used below available.  Thus regularizing the ideal
problems never requires leaving the neighborhood that preserves $\Omega$.
The ordinary library proposition itself is independent of this choice:
it applies to arbitrary ordinary data and assumes neither the metric
clamps nor a nonzero count.

Here is the analytic statement used in the dimension calculation.  Its
ordinary assertion is also available before the spinor construction and
does not assume the nonvanishing of any count.

\begin{theorem}[Compatible regularization]\label{thm:analysis:regularization}
For the fixed finite family, equation perturbations and representatives
can be chosen arbitrarily small in the prescribed norms with the following
properties.
\begin{enumerate}
\item Ordinary irreducible equations and the labelled sphere, surface,
and point tests are transverse, including the relevant ideal and broken
strata.  Their representatives have the exact support and persistence
properties of Lemma~\ref{lem:analysis:incidence}.  They are compatible with
fixed determinant gluing, marked translation on cylinders, and all true
faces.  Prescribed harmonic reducers on separated caps are preserved.
\item Single-component, single-valued equivariant perturbations give the
fixed-type regularity in Definition~\ref{def:indices:analytic-hypotheses}.
Every tuple of fixed type is consequently excluded except the legitimate
unbroken cut-down instantons.
\item A finite collection of rational multisection terms, in addition to
single-valued terms, eliminates every forbidden mixed ideal stratum by
the negative dimension of its necessary projection.  No equation or
retained cut on this projection depends on discarded $R$ fields or their
particles.  The allowed free stratum is regular of dimension one, and the
single trace-end main problem is regular of dimension zero.  Choices on
the exterior of a separated lens cap agree under the corresponding bit
flip.
\item Solutions have the compactness, removable-singularity, integral
particle-charge, and exponential true-end properties proved below.
Near each legitimate instanton, the phase-cut link has oriented count
$2^\eta$ times its instanton sign, where $\eta=n_D^{\rm tot}-1$.
\end{enumerate}
The finite branch lists have compatible weighted germs and total weight
one.  After the no-particle end gluing of Section~\ref{sec:gluing}, their
oriented one-dimensional zero sets obey the ordinary rational Stokes
rule.
\end{theorem}

We construct the representatives, prove compactness, and regularize the
fixed and forbidden strata in this section.  The proof of the theorem is
completed in Section~\ref{sec:gluing}, after the remaining lens ends have
been given regular collars.  Closed-manifold monopole compactness
provides useful background \cite{FeehanLenessPU2I}, but does not by itself
give the face, support, and projected-regularity assertions above.

The proof order is summarized in Table~\ref{tab:analysis:stages}.
Sections~\ref{sec:ordinary}--\ref{sec:stack} construct the nonzero count
$\Omega$; the ordinary library used there is independently available.
The remaining stages preserve this count and close the spinor boundary
argument.  In particular, the regular end construction in
Section~\ref{sec:gluing} precedes the final interior regularization there.
\begin{table}[htbp]
\centering
\small
\begin{tabular}{@{}p{.27\textwidth}p{.67\textwidth}@{}}
\toprule
Stage & Output used next \\
\midrule
Ordinary representatives & An independent ordinary library with exact
support rules.  Later choices remain in the neighborhood preserving
the count constructed in Sections~\ref{sec:ordinary}--\ref{sec:stack}. \\
Energy and fixed types & Ideal compactness, exclusion of unwanted fixed
types, and the legitimate instanton links. \\
Projected finite induction & Exclusion of forbidden mixed candidates
and regular zero-dimensional trace-main problems. \\
Regular end gluing & Section~\ref{sec:gluing} constructs exhaustive
collars for the remaining no-particle lens ends, preserving branch
weights and orientations. \\
Final free stratum & Section~\ref{sec:gluing} completes the interior
one-dimensional regularization and applies the weighted boundary rule
to the resulting compact count. \\
\bottomrule
\end{tabular}
\caption{Stages of the analytic construction.  The numerical
inequalities of Section~\ref{sec:indices} are applied as their
regularity and incidence hypotheses become available.}
\label{tab:analysis:stages}
\end{table}

\subsection{Function spaces and allowable supports}

Let $X$ be a completed component of an open face, and choose product
coordinates $(t,y)$ on each end.  Use a smooth reference configuration
which is the prescribed flat connection with zero spinor sufficiently
far out; exponentially convergent references give the same spaces.  Write
$c^{a,\alpha}$ for the closure of smooth sections in the $C^{a,\alpha}$
norm, with $a\geq3$ and $0<\alpha<1$.  For $\delta>0$ put
\begin{equation}\label{eq:analysis:weighted-norm}
 \|u\|_{a,\alpha,\delta}
 =\|u\|_{c^{a,\alpha}(X_{\rm core})}
  +\sup_{j\geq0}e^{\delta j}
      \|u\|_{c^{a,\alpha}([j,j+2]\times Y)},
\end{equation}
summing the last expression over the finitely many ends.  Take the
closure of compactly supported smooth sections in this norm.  We use
$E_a=c^{a,\alpha}_\delta(\Lambda^1\otimes\ad E\oplus W^+\otimes E)$
for field differences and $F_a=c^{a-1,\alpha}_\delta$ for the curvature,
Dirac, and gauge-slice equations.  On finite inserted cylinders use the
corresponding uniform strip norms; exponential weights are measured from
the ends when a gluing norm is needed.  Choose $\delta$ smaller than the
positive normal spectral gaps of the finitely many limiting flat states.
The specific asymmetric lens norm is described in
Section~\ref{sec:gluing}.

When a Sobolev cokernel test is used, choose a second weight
$0<\delta_S<\delta_H=\delta$, both below the same positive indicial
gaps.  Then $c^{a,\alpha}_{\delta_H}$ embeds in
$W^{1,p}_{\delta_S}$ for every finite $p$: on the $j$th strip the
weighted integrand is bounded by a constant times
$e^{-p(\delta_H-\delta_S)j}$.  We do not use inclusion at equal
weights, which need not hold.  Elliptic end estimates transfer
surjectivity back to weight $\delta_H$ because no indicial root lies
between these two weights, as detailed below.

The local gauge group consists of $c^{a+1,\alpha}$ determinant-one gauges
whose difference from a parallel limiting constant is in
$c^{a+1,\alpha}_\delta$.  Constants at different ends are independent.
Thus the gauge Lie algebra is the decaying space together with a finite
sum of cut-off parallel end sections.  It is not the gauge group with a
chosen frame fixed at every end.  The usual local slice is obtained by
adjoining the adjoint of the infinitesimal gauge action to the field
equations.  For a non-$R$ field its gauge stabilizer is finite.  At an $R$
field we use the Abelian equations and its parallel stabilizer instead.
If $\mathfrak s$ is its Lie algebra of global parallel stabilizers,
the augmented gauge target in this latter convention is
\[
 \operatorname{ann}(\mathfrak s)
 =\left\{f\in c^{a-1,\alpha}_\delta:
       \int_X\langle f,\sigma\rangle=0\text{ for every }
       \sigma\in\mathfrak s\right\},
\]
and gauge parameters are taken modulo $\mathfrak s$.  An augmentation
using the full target instead retains its $\mathfrak s$ cokernel,
as on the cap in Section~\ref{sec:gluing}.  All universal
surjectivity assertions below concern non-$R$ fields; in particular
the nonzero active spinor of type $S$ kills its infinitesimal
determinant-one gauge stabilizer.
The weighted elliptic operator on each fixed-asymptote slice is Fredholm;
the extra limiting constants implement precisely the index convention
of Section~\ref{sec:ordinary}.  The fixed-operator weighted Sobolev
Fredholm theory is classical
\cite[Theorems~6.2 and~8.1, Lemmas~7.1 and~7.3]{LockhartMcOwen1985}.
The transfer to the present little-H\"older spaces and the estimates
uniform in neck length are established below.  The local analytic
inputs are the small-curvature gauge theorem
\cite[Theorems~1.3 and~2.1, Corollary~2.2]{Uhlenbeck1982Connections},
the usual gauge slices~\cite{DonaldsonKronheimer1990}, and local
elliptic parametrices and Sobolev/Schauder estimates
\cite[Corollary~2.1.B and Theorem~2.2.C, revised author text]{Taylor1991}.

The local sampling terms used below read holonomies and transported
spinor values on compact graphs of loops and paths, and return outputs
supported in compact balls.  We first specify how these supports behave
when necks separate.

The parameter space is treated stratum by stratum.  Near a corner its
length coordinates $T_j$ are independent, and the embeddings of each
compact subset of a limiting component commute.  A datum on a face is
extended using these embeddings and a product of parameter cutoffs.
A connected sampling graph and every ball in its energy cutoff belong
to the same component.  A term requiring transport through a cut is
turned off before that cut separates.  A nonstationary term far down an
inserted collar is either zero there or has an exponentially small
coefficient, with the same statement for its derivatives.  We may impose
any fixed finite exponential rate required subsequently.  This is a
condition on terms extending to increasing depths, not a restriction on
compactly supported tests on any given open face.

Here is an explicit form of the exponential requirement.  For an
elementary term $p$, let $p^{[D]}$ mean its contribution whenever an
input or output graph reaches depth at least $D$ in an inserted
collar, or crosses a neck of length at least $2D$; include any
nonstationary error in this contribution.  Put
\begin{equation}\label{eq:analysis:exponential-seminorm}
 N_{k,\mu}(p)=\sup_{D\geq0}e^{\mu D}
       \|p^{[D]}\|_{C^k_{\rm par,field,space}},\qquad \mu>0,
\end{equation}
on each of the bounded configuration and parameter sets used for
the library norms.  The norm includes the relevant output operator
derivatives, and its supremum is uniform in the other length
coordinates.  Each term has bounded-depth compact graphs on its
limiting part, and its incompatible crossing extension is cut off
at a finite length, so these seminorms are finite.  Include integral
$k,\mu$ in the countable norm list below.  The summed contributions
affected by flattening a middle of depth $D$ are then
$O(e^{-\mu D})$ at every fixed required order.  Fixed polynomial
factors from derivatives can be absorbed by decreasing $\mu$.
Parallel end constants do not supply a nondecaying forcing term:
their temporal part is removed by the parallel-group gauge ODE,
and its endpoint rotations are allowed independently.

There are two further restrictions.  First, a single-component term
between two $J$ cuts uses only the metric parameters retained on that
component: those in the intervals from its left bounding interval to its
right bounding interval, inclusive.  It vanishes near any intermediate
$J$ separation incompatible with that component.  Second, terms reaching
a separated lens cap $N$ never couple it to another component.  The
Dirac perturbation on $N$ is zero, the curvature perturbation is zero on
its Abelian locus, and exterior data are equal under the bit flip.
Inside $N$ use the tensor-equivalent ASD data and sphere tests.  These
requirements are simultaneous at corners, by the product collars of
Theorem~\ref{thm:geometry:family}.  They leave all compact probes on the
relevant main components available.

Ordinary perturbations on a finite $Y$ seam are time-independent gradients
of smooth invariant functions of simultaneous solid-torus holonomies,
averaged with a transverse disk density.  The thickenings are submersive,
their values are bounded, and their spatial Hessians are bounded on
$L^2$.  Retain their ordinary transversality choices and extend them
independently of the spinor.  They are supported away from the true
cylinders.  Their gradient nature, rather than a bound integrated over
the seam length, is what is used in the energy estimate of
Lemma~\ref{lem:estimates:cylinder}.
For the spinor family a fixed finite $Y$ seam is part of the compact
body.  Additional equation probes may be supported in any of its
compact balls; their norms are included in the small nongradient
remainder after the seam lengths have been fixed.  This permits full
local output variation there.  The equation on a genuine ordinary
Floer cylinder, where the length is unbounded in Floer theory, retains
the separate time-independent gradient prescription.

\subsection{Exact representatives}

\begin{lemma}[The zero-winding modification]\label{lem:analysis:zero-winding}
Fix $1/2<b<1$.  On the Hilbert manifold of $H^b$ paths
$g:[0,1]\longrightarrow\SU(2)$ with endpoints $1$ there is a smooth
conjugation-equivariant operation $g\mapsto\widehat g$, homotopic to the
identity operation relative to the endpoints, such that every commuting
path of winding zero has image a uniform positive distance from $-1$.
The operation fixes commuting paths of nonzero winding.  It uses only
the given path.
\end{lemma}
\begin{proof}
Identify $\SU(2)$ with the unit quaternions.  For a nonconstant commuting
path define
\[
 M(g)=\int_0^1\operatorname{Im}g(t)\otimes
                    \operatorname{Im}g(t)\,dt.
\]
This is a rank-one positive matrix: if its eigenvalue were zero, continuity
would force $g$ to be the constant path $1$.  Near such a path its largest
eigenvalue remains simple and determines a smooth, unoriented real line
$L(g)$.  Orthogonal projection onto $\R\oplus L(g)$ followed by
normalization gives a path $h_g$ in the circle subgroup associated to
$L(g)$.  Shrinking the neighborhood makes this projection nonzero at
every time.  Sobolev embedding $H^b\hookrightarrow C^0$ and the Banach
algebra property for $b>1/2$ show that all these operations are smooth.

In the winding-zero neighborhood, the lift of $h_g$ starting at zero is
a unique $H^b$ path $s_g(t)\in L(g)$ with $s_g(0)=s_g(1)=0$ and
$h_g=\exp s_g$.  This definition is independent of an orientation of
$L(g)$.  Let $C$ be the set of commuting winding-zero paths meeting
$\{\operatorname{Re}\leq0\}$.  It is closed in the full based path space.
Indeed axes have a convergent subsequence, uniform convergence preserves
commutativity, and a limiting path meeting this hemisphere is
nonconstant; the degree of its projected circle path is locally
constant.  An invariant open neighborhood $U$ of $C$ can consequently
be chosen on which $L(g)$ and $s_g$ are defined and which contains no
commuting path of nonzero winding.

Smooth separation on a Hilbert manifold gives a function $f$, equal to
one on $C$, with support in $U$ and $0\leq f\leq1$.  Averaging over
conjugation preserves these properties.  Set
\[
 \widehat g(t)=g(t)\exp(-f(g)s_g(t)),
\]
with the second factor equal to one outside $U$.  Multiplying the exponent
by a parameter in $[0,1]$ supplies the required homotopy.  On $C$ the
modified path is constant.  A commuting winding-zero path outside $C$
lies in the positive-real hemisphere; its lifted angles lie in
$(-\pi/2,\pi/2)$ and partial multiplication replaces them by
$(1-f)$ times those angles.  It remains in that hemisphere.  Thus all
such modified paths lie in the closed positive-real hemisphere, whose
distance from $-1$ is positive.  The stated behavior at nonzero winding
follows from the choice of $U$.
\end{proof}

Apply the lemma to the normalized holonomy sweep of each sphere map.
Parallel transport compares the frames, and the continuous square root
of the determinant holonomy, starting at one, makes the sweep
$\SU(2)$-valued.  The determinant evaluations on these spheres are even,
so the sweep ends at one.  At a parallel splitting its eigen-winding is
$v(S)/2$.  The equation $\widehat g(t)=-1$ has real codimension three with
one moving mark $t$, hence degree two.  Subsequent equivariant cut
variations are absolutely smaller than half the avoidance gap in
Lemma~\ref{lem:analysis:zero-winding}.  Include all comparison paths in
the basic support.  The support is the image of a compact space of
dimension at most two, and no thickening is added to it for purposes of
incidence.

\begin{lemma}[Segmentation at a primary cut]\label{lem:analysis:segmentation}
The sphere representatives can be chosen so that, at a true $J$ face,
their event is the disjoint union of the events on the two half spheres,
with the prescribed signs.  During the change from the whole-sphere rule
to this rule, a commuting path whose two half degrees vanish retains a
uniform avoidance gap.  Every event on the face is intrinsic to its
assigned side; no derivative of transport along the growing neck enters
its definition.
\end{lemma}
\begin{proof}
The punctured traces are simply connected.  The Hurewicz identification
of $\pi_2$ with $H_2$ shows that the sphere class and the signed sum of
the two half classes are represented by homotopic sphere maps in the
interval under consideration.  Realize the latter map by concatenated
sweeps with identity intervals and fixed whiskers.  Use the whole-path
operation until concatenation has been reached.

At concatenation write $s$ for the whole logarithm and $s_1,s_2$ for
the logarithms of the two zero-degree segments, extended by zero on
the other segment.  On commuting data with both degrees zero,
$s=s_1+s_2$.  Interpolate the multipliers by
\[
 \exp(-(1-u)f s)\,
 \exp(-u f_1s_1)\,\exp(-u f_2s_2),\qquad 0\leq u\leq1.
\]
At a point of a segment meeting the nonpositive hemisphere its own cutoff
is one, as is the whole cutoff.  The exponent therefore cancels its
whole angle.  If a segment remains in the positive hemisphere, any
partial cancellation still keeps its angle between its original value
and zero.  This proves the same uniform gap throughout the interpolation.
The interpolation is equivariant and is defined on neighborhoods by the
same cutoffs as before.  The metric construction places this switch in
the regimes where the old whole-sphere clamp is no longer required on
that side.

After switching, keep marks off the identity intervals.  A whisker to
the second sweep and its inverse conjugate that sweep by a transport
independent of its internal mark.  The equation that its value is $-1$
and the modification are conjugation-invariant, so they may instead be
defined in a frame intrinsic to that half.  Variations of its cut use
the same intrinsic description.  Whiskers and the earlier homotopy lie
in the original interval interior.  At the actual split there is thus
no graph crossing $J$ and no length-dependent frame comparison.
\end{proof}

To retain support incidences through finite switches, retain the earlier compact
homotopy-support domains on their closed finite parameter ranges and
the intrinsic half-support domains on their subsequent closed ranges.
At a switching endpoint use the union of their limiting supports.
This is a finite collection of closed incidence charts of dimension
at most two in each parameter fiber; transversality is imposed also
on their boundary strata, whose parameter dimension can only decrease.
Consequently existing hits persist across every finite switch.
In a neighborhood of the true face only the later intrinsic half
domains remain.  This convention changes neither the smoothly
switching event equation nor the boundary of the actual zero set.
It does not thicken a support or increase the released-degree budget.

For a cap surface or point use two complex sections of the rank-three
universal adjoint bundle at the evaluation point, and require the
resulting $3\times2$ matrix to have rank at most one.  Its rank-one
stratum has complex codimension two; the rank-zero stratum has complex
codimension six.  This is the usual degeneracy representative of
$c_2(\ad_{\C}E)=-p_1(\ad E)$, with the sign and normalization for the
chosen insertion kept outside the count.  A variable point on an
immersed surface reduces the real degree from four to two.  Choose
distinct surface representatives pairwise transverse with no triple
incidence; fixed cap points are general.  All sphere supports are
disjoint from one another and from the cap tests.  This is the exact
evaluation representative: the sampling graphs used to construct its
sections do not enlarge its basic incidence support.

Phase cuts are sections of character two for the original spinor phase
$\mathbb T$.  They are determinant-one gauge invariant.  Equation
perturbations use none of the sweep marks or surface evaluation variables.
A cut variation uses at most its own such variable, and phase sections
use none.  These restrictions ensure the index accounting in
Section~\ref{sec:indices}.

\subsection{Sampling libraries and the finite isotropy construction}

For each connected graph choose one base frame.  Its data consist of
finitely many based holonomies and spinor values transported to that
frame.  Use bounded spinor samples, for example
$\Phi/(1+|\Phi|^2)^{1/2}$, when a global bound is needed.  An output is
a fixed smooth form or spinor supported in a ball, transported from the
base, times a smooth bounded equivariant function of these data.  The
same construction gives section values and small group-valued variations
by exponentiation.  Include ordinary compact cutoffs in sample and
parameter charts.

For a finite small-ball cover of every input graph and output region
choose functions $\chi_\nu$ which are one on smaller balls still covering
these sets.  Multiply the term by
\begin{equation}\label{eq:analysis:energy-cutoff}
 \prod_\nu\vartheta_\nu\left(
  \int_X\chi_\nu\bigl(|F_a^0|^2+|\Phi|^4\bigr)
          +8\pi^2\sum_j m_j\chi_\nu(x_j)\right),
\end{equation}
where $\sum_jm_j[x_j]$ is the current ideal label, and each
$\vartheta_\nu$ is supported below a small regular-gauge threshold.
Its value can be one on the smaller neighborhood being tested.  The
threshold is also smaller than a charge-one curvature atom.  A finite
ball cover can always be made sufficiently fine at a smooth field and
away from specified particles.  Crucially, each $\chi_\nu$ is confined
to the component of its associated graph.  Weighted atoms enter this
formula additively, including at collisions.

Enumerate a dense countable library of these single-valued terms by
$p_j$.  Choose positive weights $w_j$ as follows.  List the desired
fixed-order operator bounds on bounded configuration and parameter sets,
the corresponding active-graph regularity bounds on nested balls, and
the exponential end bounds, as a countable list of seminorms $N_k$.
Arrange
\begin{equation}\label{eq:analysis:weights}
 w_j\geq 2^j\max(1,N_1(p_j),\ldots,N_j(p_j)),
 \qquad
 \|\lambda\|_{\mathfrak P}=\sum_j w_j|\lambda_j|.
\end{equation}
For each fixed $k$ the finitely many ratios $N_k(p_j)/w_j$ with $j<k$
have a finite maximum $C_k$; for $j\geq k$ the ratio is at most $2^{-j}$.
Thus $\sum_j\lambda_jp_j$ converges in every listed norm, with a bound
depending on $k$, uniformly on bounded sets, and tails tend to zero there.
The parameter space $\mathfrak P$ is a separable weighted $\ell^1$ space.
Taking a small ball in it makes any prescribed finite collection of
equation norms as small as required.  The field derivative of each term is compact from $E_a$
to $F_a$: graph evaluation and transport are of order zero, inclusion
from field order $a$ to equation order $a-1$ is compact on its compact
support, and derivatives of the energy cutoffs are finite-rank scalar
functionals.  The tail estimates preserve compactness for the sum.

\begin{lemma}[Finite samples and available variations]
\label{lem:analysis:samples}
Finite framed samples, avoiding any specified finite particle set,
detect the stabilizer of a smooth field.  They give arbitrary curvature
and retained cut variations at type $A$, and arbitrary tangential
curvature and active-line spinor variations at type $S$.  With a type
$P$ witness they give arbitrary full equation and retained cut variations
at every non-$R$ target, using finite weighted branch lists if needed.
Phase-cut values can be varied using the witness alone.
\end{lemma}
\begin{proof}
A gauge stabilizing a connection is parallel and is determined by its
base value, which commutes with all based holonomies.  It must also fix
all transported spinor values, up to the prescribed common phase when
phase is included.  For $\SU(2)$ one noncentral holonomy reduces its
centralizer to a circle and two noncommuting ones reduce it to the
center.  Spinor values impose finite-dimensional linear eigenspace
conditions.  This proves finite detection; equivalently the stabilizer
is the common stabilizer of finitely many matrices and vectors in these
fixed finite-dimensional representations.  Removing finitely many
points does not change the parallel-transport criterion.  Paths can be
perturbed to avoid these points, retaining the finitely many strict
conditions that detect the stabilizer.

At $A$ the frame stabilizer is the center, which acts trivially on the
adjoint target.  At $S$ the combined frame and phase stabilizer preserves
the ordered active line.  Its invariant outputs are exactly diagonal
curvature and spinors in that line.  In a slice tube for the compact
sample group an arbitrary invariant output at the center extends to an
equivariant function: first choose it on the slice, average over its
stabilizer, and extend by the group action.  Cutoffs supported inside
the tube globalize it.  Outputs in arbitrarily small balls span all
smooth compactly supported directions in the stated target.  Separate
cut coefficients give their values independently of equation outputs.
The sample and parameter cutoffs may be supported away from the
zero-spinor reducible locus.  At a central end, choices of a parallel
line create no extra parameter after unframing: its full connected
parallel group acts transitively on these lines and its elements extend
over the collar.  The residual group is the splitting-preserving group
already used for the tangential equations.

For the mixed assertion use independent frames for the target and the
witness, together with the common phase.  If a parallel gauge on a
$P$ field compensates a phase $z\ne\pm1$, its distinct eigenspaces are
parallel lines and the spinor lies in one of them.  That would be type
$S$.  Thus the projection of its stabilizer to phase is exactly
$\{1,-1\}$.  With phase so restricted an $A$ target has only its central
stabilizer and an $S$ target only a finite stabilizer.  The product sample
group therefore has finite isotropy $H$.

In a tube $G\times_HD$, choose an output $s$ on the slice and include
all translates $h\cdot s(h^{-1}\cdot)$, $h\in H$, each with weight
$1/|H|$.  Multiply by an invariant cutoff whose support lies inside
the tube.  On a neighborhood of its boundary every branch is zero;
extend outside by $|H|$ repeated zero branches.  A change of slice
representative permutes the list with multiplicities, so this is a
well-defined weighted germ.  Choose the coefficient space to contain
the whole $H$-span of each chosen output space.  Every branch then has
the same full coefficient range at the center, giving surjectivity
branch by branch.  A witness alone has stabilizer acting trivially on
the phase-character-two target, so arbitrary phase-section values are
available there even single-valuedly.
\end{proof}

\begin{proposition}[Ordinary exact-support library]
\label{prop:analysis:ordinary-library}
The preceding construction at zero spinor supplies the ordinary
representatives and perturbations used in Sections~\ref{sec:ordinary}
and \ref{sec:caps}.  It is independent of the existence or value of the
closed cap count.  It permits simultaneous regularity on the countably
many required lower strata and normal derivative regularization at the
specified harmonic reducers while fixing their values.
\end{proposition}
\begin{proof}
At an irreducible instanton all required adjoint outputs are available
by Lemma~\ref{lem:analysis:samples}.  Universal surjectivity follows by
pairing a cokernel with arbitrary compactly supported smooth outputs
away from the finitely many labelled points; these outputs are dense
in $L^p$.  Independent variations give the full value matrix of each
retained rank test and the three holonomy-event directions.  The
transversality argument below applies on each position, rank, parameter,
and topology stratum.  It uses no nonvanishing input.

For a specified harmonic reducer take a sample slice with its circle
stabilizer singled out by noncentral data.  Infinitesimal based
holonomies detect a connection variation modulo infinitesimal gauge:
if every loop variation is simultaneous infinitesimal conjugation,
transport that infinitesimal base value along paths; the loop identity
makes the result independent of the path and its covariant derivative
is the given variation.  A finite number of samples thus injects the
finite-dimensional normal kernel.  Localized output forms detect its
finite-dimensional cokernel.  Slice-linear equivariant variations
therefore prescribe any complex-linear kernel-to-cokernel map in the
normal circle representation.  They vanish on the commuting slice and
leave the Abelian equation unchanged.  A generic arbitrarily small map
makes a complex index-zero block invertible, or a complex index-one
block surjective.  Independent cut variations prescribe the corresponding
normal cut derivative.  Subsequent irreducible variations are supported
off the reducible locus.

On an ordinary cylinder use the time-independent gradient library for
the equation, whose irreducible transversality is established in
Section~\ref{sec:ordinary}.  For a cut based at a moving mark translate
its graph with that mark.  This is joint translation invariance of the
marked problem; the equation still does not depend on the mark.
Energy cutoffs and labelled atoms translate with it.  The face and end
prescriptions above make these choices agree with compact probes on
the pieces and with their exponentially small extensions.  The gauge
group and the fixed transition identifications are those already
specified, so this construction introduces neither an $\SO(3)$ lift
sum nor a new matching factor.  Persistence and particle accounting
are proved in Lemma~\ref{lem:analysis:incidence}.
\end{proof}

\begin{lemma}[Independence from discarded fields]
\label{lem:analysis:R-independence}
The coupled terms can be supported so that they vanish on every tuple
of fixed type.  On an exact $J$ face a term involving a connected graph
in an $R$ component vanishes on a neighborhood of that sample.  After
discarding $R$ fields, their particles, and their tests, every remaining
equation and cut is a function only of the retained fields and particles
and of the retained parameter variables.  The same assertion holds for
separated lens-cap data.
\end{lemma}
\begin{proof}
Support a mixed term compactly inside the finite-isotropy locus of its
independent-frame sample space.  In a tuple of fixed type the common
phase can be compensated independently on every $S$ component, while
it acts trivially on every zero-spinor component.  The combined
stabilizer is therefore continuous.  An $R$ input or output has a
continuous independent frame stabilizer even when another component
contains a witness.  Both types of sample are outside the closed
support of the term.

No frame is compared between the disconnected graphs.  On a split
face the only nonzero coupled terms consequently sample retained
components.  The energy cutoffs of those terms have balls on precisely
those components, so the atomic summands in
\eqref{eq:analysis:energy-cutoff} contain no discarded positions.
First-stage terms are single-component and already have this property.
Parameter cutoffs use only allowed variables; the parameters of an
$R$ component may remain in the projection, as stipulated in its
dimension count.  Lens-cap sampling and coupling have been disabled
separately.  This proves independence for the entire formula, including
the scalar cutoff factors, not just for its visible output.
\end{proof}

\subsection{Differentiability and universal regularity}

Moving a sampling path is not a smooth operation at unchanged field
regularity.  We use the following precise version of the regularity
argument, rather than differentiating translations indefinitely on one
Banach space.

\begin{lemma}[Compatible implicit charts]\label{lem:analysis:scale}
Let $E_s=c^{a+s,\alpha}_\delta$ and
$F_s=c^{a+s-1,\alpha}_\delta$, with the finite-dimensional equation,
parameter, and cut variables included.  For the sampling formulas above,
the equation map at level $s$ is $C^1$.  For every prescribed $j$ it is
$C^j$ from sufficiently high field order to $F_s$.  At a smooth
universally surjective zero its universal solution set has $C^j$ local
charts, for every prescribed finite $j$, with the same perturbation
space $\mathfrak P$ and the same smooth finite-dimensional kernel
coordinates at all levels.
\end{lemma}
\begin{proof}
The derivative of evaluation along a moving path contains one derivative
of the field; transport satisfies the same assertion by differentiating
its ordinary differential equation.  Little-H\"older translation is
continuous, so field order $a+s$ gives a continuous first derivative
with target order $a+s-1$.  Additional parameter differentiations use
additional field derivatives.  Smooth functions, finite products,
integration in \eqref{eq:analysis:energy-cutoff}, and the summability
in \eqref{eq:analysis:weights} preserve these assertions.  Metrics and
bundles are compared in fixed local identifications, so their variation
only changes smooth coefficients of the differential operators.

Write the field linearization as $L_s:E_s\to F_s$.  It is elliptic plus
a compact derivative.  At a smooth solution its kernel is smooth and
independent of $s$, and higher-order right-hand sides give higher-order
solutions.  For the graph terms this is the same elliptic bootstrap as
below: their field derivatives have order zero and the cutoff derivatives
are scalar functionals multiplying smooth outputs.  Pick finitely many
smooth perturbation directions $Q\subset\mathfrak P$ spanning the
cokernel at the lowest level.  Linear regularity shows that
$\widetilde L_s:E_s\oplus Q\to F_s$ is surjective at every level and
has a common smooth finite-dimensional kernel $K$.  Choose a projection
onto $K$ using finitely many compactly supported smooth pairings and
coordinates on $Q$.  Its kernels $V_s$ are compatible complements, and
$\widetilde L_s:V_s\to F_s$ is an isomorphism.

Write $\mathfrak P=Q\oplus\mathfrak P'$ and
$Z=K\oplus\mathfrak P'$.  The $C^1$ implicit theorem at each level
gives a graph $v_s:U_s\subset Z\to V_s$.  On sufficiently small
common domains uniqueness of the lowest-level graph makes the graphs
agree.  For any finite list of higher levels we may intersect their
domains; thus the common graph $v$ is $C^1$ into every level needed in
a given finite-order argument.

Here is the higher-order step.  In graph coordinates write
$\mathcal F(v(z),z)=0$ and
$B_s(z)=D_v\mathcal F(v(z),z):V_s\to F_s$.  It remains invertible
near the base point.  If $v$ is $C^r$ into the requisite higher levels,
its differentiated equation has the form
\[
 B_s(z)H_r(z)=R_r(z),\qquad H_r=D^rv,
\]
where $R_1=-D_z\mathcal F$, and for $r\geq2$ the right-hand side
uses derivatives of $v$ only up to order $r-1$.  Choose a level $t$
high enough that $\mathcal F:V_t\times Z\to F_s$ is $C^{r+1}$.
Then $R_r$ and $B_s|_{V_t}$ are $C^1$, while $H_r$ is continuous
into $V_t$.  Subtracting the equations at $z+h$ and $z$ gives
\begin{align*}
 B_s(z)\bigl(H_r(z+h)-H_r(z)\bigr)
  ={}&R_r(z+h)-R_r(z)\\
    &-\bigl(B_s(z+h)-B_s(z)\bigr)H_r(z+h).
\end{align*}
Apply the fixed inverse $B_s(z)^{-1}$.  The right-hand side has a
first-order expansion with $o(\|h\|)$ remainder, using continuity of
$H_r$ at level $t$.  This proves differentiability into $V_s$; the
same expression proves continuity of the derivative.  Induction gives
$C^{r+1}$ charts at the lower level.  Only finitely many higher charts
are needed for any specified order.  In particular this proof does not
differentiate the inverse repeatedly as a same-order operator.
\end{proof}

To test universal surjectivity, first use the augmented linearization
$W^{1,p}_{\delta_S}\to L^p_{\delta_S}$ with $p>4$ at a fixed smooth
zero in the weight-$\delta_H$ configuration space.
The perturbation outputs span the curvature, Dirac, and cut targets,
not the added gauge-adjoint target.  The latter cokernel component is
removed first: pairing with infinitesimal gauge variations gives the
gauge Laplacian $G^*G$, where $G$ is the infinitesimal gauge action.
A dual solution $\xi$ has weight $-\delta_S$.  On an end the
limiting gauge operator is $-\partial_t^2+\Delta_\gamma$.
Its spectral decomposition and elliptic regularity give
\[
 \xi(t)=a_e+t b_e+O(e^{-\epsilon t}),\qquad
 a_e,b_e\in\ker\Delta_\gamma,
\]
with exponential decay in the complementary modes.  Growing normal
modes are excluded by the choice of weight.  Exponentially decreasing
coefficient errors preserve this expansion, after decreasing
$\epsilon$ to absorb polynomial factors.  For every independently
allowed cut-off parallel end constant $\chi_e c$, pair the cokernel
with the field variation $G(\chi_e c)$.  Gauge equivariance kills
the equation and cut rows.  Green's formula for the gauge row gives
\[
 0=\int_X\langle\xi,G^*G(\chi_e c)\rangle
   =-\langle b_e,c\rangle_{L^2(Y_e)}.
\]
Thus every $b_e$ is zero.  The dual solution is bounded with decaying
derivative, so a second integration by parts gives
$\|G\xi\|_2^2=0$.  It is a global infinitesimal stabilizer and is
zero on the non-$R$ slices used here.  More generally the extended
gauge Laplacian has kernel and cokernel $\mathfrak s$, and induces
an isomorphism from the gauge domain modulo $\mathfrak s$ to
$\operatorname{ann}(\mathfrak s)$.  This also explains the target
convention specified above.  The argument applies to tangential
non-$R$ slices and to infinitesimal local lifts of finite branches.

Compactly supported smooth outputs away from finitely many particles
are dense in its equation target.  The variations in
Lemma~\ref{lem:analysis:samples} consequently annihilate no nonzero
cokernel in the required full or tangential target.  Field evaluation
on the compact paths is continuous at this level.  Independent section
variations give the retained finite-dimensional cut targets.

Linear elliptic regularity gives the same surjectivity at higher orders and
weight $\delta_H$.  To justify the weight step explicitly, an
$L^p_{\delta_S}$ solution with right-hand side in
$c^{a-1,\alpha}_{\delta_H}$ satisfies on each sufficiently remote
end the limiting spectral equation plus a small decreasing-coefficient
error and the summable graph tails.  The positive and negative spectral
integrals improve its decay to $\delta_H$, since the interval
$[\delta_S,\delta_H]$ contains no indicial root; parallel gauge
modes have just been treated by their separate end constants.
The error is absorbed in the weight-$\delta_H$ norm, and graph tails
in \eqref{eq:analysis:exponential-seminorm} are chosen faster than
$\delta_H$.  Interior Schauder estimates then give the asserted
H\"older order.  The same argument identifies the smooth kernels and
the finite-dimensional cokernel quotients at both weights, including
the finite extended gauge coordinates.  The cokernel identification
uses the range and regularity of solutions, not smooth representatives
of the full adjoint cokernel: adjoints of thin samples can have path-
or point-supported distribution terms.  Density of smooth outputs
detects those cokernels without any smoothness assertion.
It does not presume that equal-weight H\"older
closure embeds in weighted Sobolev space.

Equivariance makes gauge directions redundant, also on local branch
lifts.  Apply parametric transversality in the charts of
Lemma~\ref{lem:analysis:scale}; the differentiability order can always
be chosen larger than both zero and the relevant Fredholm index
\cite[Theorem~1.3 and Corollary~1.5]{Smale1965}.
Finite isotropy is handled on the local finite slice
covers.  Rank-one and rank-zero matrix strata are treated separately,
with variations of the full matrices available before either restriction.

\subsection{Energy, regularity, and ideal limits}

\begin{lemma}[Compactness on bodies and removable particles]
\label{lem:analysis:compactness}
For the fixed finite family, and for perturbations in a sufficiently
small bounded set of the norms above, the following statements hold
uniformly, including coefficient-convergent sequences of already ideal
solutions.
\begin{enumerate}
\item The trace-free curvature and covariant spinor gradient are bounded
in $L^2$, and the spinor is bounded in $L^4$.
\item The measure $|\Phi|^4\,d\Vol$ has no atoms in a weak limit,
including a limit obtained by translating along a true cylinder.
Neither does $|F_a^{0,+}|^2\,d\Vol$.
\item On compact sets, a subsequence converges modulo gauge away from
finitely many points.  The limiting field extends across these points
and solves the same equations with the ideal cutoff formula
\eqref{eq:analysis:energy-cutoff}.  The lost charges are positive
integers, and each charge-$m$ atom has mass $8\pi^2m$ in the stated
curvature norm.
\item On any fixed ideal stratum, if there is no further bubbling or
end loss, convergence is in the completed higher-order field charts.
At a new bubble there is still convergence in $C^{0,\beta}$ on surviving
exact cut paths for any fixed $\beta<1$.
\end{enumerate}
\end{lemma}
\begin{proof}
The integrated inequality of Lemma~\ref{lem:estimates:local}, combined
with Proposition~\ref{prop:estimates:energy} and Chern--Weil, gives
\begin{equation}\label{eq:analysis:global-energy}
 \|F_a^0\|_2^2+\|\nabla_a\Phi\|_2^2+\|\Phi\|_4^4\leq C.
\end{equation}
Here the negative zeroth-order coefficients have uniformly bounded
squared integral.  The true collar middles have positive scalar
curvature, asymptotically flat background, and exponentially small
forcing; the remaining finite body, including all finite $Y$ seams,
has a fixed bound.  The perturbation values have bounded squared
integral by construction.  On completed ends first use decaying cutoffs
and then pass to the limit.  The integrated Dirac identity uses
$\|P_D\|_2^2$, not a derivative of $P_D$.  The curvature identity is
\[
 \|F_a^0\|_2^2=2\|F_a^{0,+}\|_2^2+8\pi^2\kappa.
\]
The charge bound and the positive quartic term therefore give all
three terms in \eqref{eq:analysis:global-energy}.

For the nonconcentration statement take a cutoff equal to one on $B_r$
and supported in $B_{2r}$.  The local inequality gives, with harmless
positive constants and uniformly bounded local geometry,
\begin{equation}\label{eq:analysis:no-atom}
 c\int_{B_r}|\Phi|^4
 \leq C r^{-2}\int_{B_{2r}\setminus B_r}|\Phi|^2+o_r(1)
 \leq C\left(\int_{B_{2r}\setminus B_r}|\Phi|^4\right)^{1/2}
          +o_r(1).
\end{equation}
The lower-order and bounded inhomogeneous terms tend to zero on these
shrinking balls; a uniform $L^4$ bound and H\"older's inequality give
this also for a bounded coefficient times $|\Phi|^2$.  First pass to
a weak measure limit at radii with zero boundary mass, then let $r$
tend to zero along such radii.  The limiting annulus masses tend to
zero even if there were an atom at the center.  Equation
\eqref{eq:analysis:no-atom} forces that atom to vanish.  The same local
argument applies after translation on a true cylinder.  Since
$F_a^{0,+}$ is quadratic in $\Phi$ plus a bounded perturbation,
its squared norm has no atom either.

Fix a ball outside curvature concentration points.  Small-energy
Coulomb gauge gives a $W^{1,2}$ connection form $a$ small in $L^4$
after choosing the ball and the energy threshold small enough.  The
spinor has $W^{1,2}$ control from
\eqref{eq:analysis:global-energy} and is small in $L^4$ on this ball,
by the nonconcentration just proved.  In these gauges the augmented
first-order system for $u=(a,\Phi)$ is
\begin{equation}\label{eq:analysis:critical-system}
 \mathcal D u=C(u)u+b,
 \qquad |C(u)|\leq C|u|,
\end{equation}
where $b$ is uniformly bounded; bounded lower-order linear terms may
be included in $\mathcal D$ or $b$.  For $2<p<4$, the critical Sobolev
estimate gives on a smaller ball
\[
 \|u\|_{W^{1,p}}
 \leq C_p\bigl(\|u\|_{L^p}+\|b\|_{L^p}
                    +\|C(u)\|_4\|u\|_{4p/(4-p)}\bigr).
\]
Absorb the last term by making the $L^4$ coefficient small.  For a
weak field the same step is justified by a local parametrix: freeze
the principal symbol on a sufficiently small chart, absorb its small
variation and the $L^4$ potential, and solve at exponents $2$ and $p$.
The cutoff commutator contains $u\in L^4$, hence belongs to $L^p$;
uniqueness of the small-potential inverse identifies the two solutions.

Choose $p>8/3$, so $4p/(4-p)>8$.  The quadratic terms are then in an
$L^q$ space with $q>4$.  A second elliptic estimate gives local
boundedness and a positive H\"older exponent.  Equation
\eqref{eq:analysis:critical-system} now gives $W^{1,q}$ for every
finite $q$, hence $C^{0,\beta}$ for every fixed $\beta<1$ on smaller
balls.  This part of the argument needs only bounds on perturbation
values, so it precedes any differentiation of their graph formulas.

For an active graph, its cutoff balls have this regularity with uniform
bounds and the graph lies in their smaller interiors.  Transport and
sample formulas consequently acquire their first regularity bounds.
When there is no ordinary $Y$ term sampling a developing concentration
point, differentiating the equations and using nested smaller radii gives
the next elliptic bounds.  Repeating this proves every fixed order.  The
diagonal choice \eqref{eq:analysis:weights} includes the constants at
each step, so the countable sum obeys the same bounds.  Transition
gauges have corresponding bounds from the connection transformation
equation.  This gives full bounds on a no-new-bubbling stratum and
smoothness of a single weak solution once it extends over the labelled
points.  It does not assert uniform high derivatives for an ordinary
$Y$ term whose loop family passes near a new particle.

There is one separate convergence issue for the ordinary solid-torus
$Y$ terms.  Thicken every transport leg, including any comparison
whisker, with the transverse disk of bases; no fixed common whisker
or common point is left in the family.  On each leg the map from
its three-dimensional path-and-disk domain to the spatial three-manifold
is submersive in its interior.  The preimage of a spatial particle
is discrete there, so its projection to the disk is a null set;
the moving endpoint maps give the same conclusion on boundary pieces.
There are only finitely many legs in a term.  Alternatively, for
these slice-wise gradient terms, all input paths avoid a finite
four-dimensional particle set at every time except its finitely many
time coordinates, which already suffices for almost-everywhere
spacetime convergence.  On the good loops transport converges, and
the bounded integrands give strong local $L^p$
convergence of the averaged perturbation values by dominated
convergence.  Thus the equations and local elliptic estimates give
strong first-order convergence off the particles.  Interpolation with
the uniform higher-integrability bounds gives $C^{0,\beta}$ convergence
on compact good sets for any prescribed $\beta<1$, in particular with
$\beta>b$ as needed for the modified sweeps.  If no new bubbling
occurs, the full graph bootstrap, including these thickened-loop terms,
gives convergence at every order.  We require only the stated lower
order convergence at a new bubble.

Choose the concentration threshold before the graph cutoff thresholds.
Every new concentration point carries at least that fixed positive
amount of curvature energy.  A graph meeting it has a covering cutoff
which equals one there and is therefore disabled.  For any other graph,
weak convergence of the energy measures gives exactly the scalar
cutoff limit in \eqref{eq:analysis:energy-cutoff}.  This proves
convergence of the limiting equations wherever they are defined;
no transport or evaluation at an atom is required.  The energy bound
leaves only finitely many new points.

For completeness, the limiting punctures are removable in the required
perturbed problem.  On a sequence of shrinking spherical annuli the
limiting $L^2$ curvature tends to zero.  Rescaled annular Coulomb gauges
are then small in $W^{1,2}$ and $L^4$; the flat annular limit is trivial
since the annulus retracts to $S^3$.  Cutting the connection to zero
inside such an annulus costs arbitrarily small $L^2$ curvature: the
cutoff derivative is controlled by the rescaled $L^4$ bound.  Apply
small-ball gauges to these fillings and pass weakly to obtain a
$W^{1,2}$ extension in an extended bundle.  On the fixed outer collar
the gauges give the same original connection up to a bundle
identification.  The spinor extends in $W^{1,2}$, and point cutoffs show
that the equations extend distributionally.  Terms whose graphs meet
the labelled atom are already off.  The preceding weak bootstrap makes
the extension regular.  This is the local small-energy removal
argument of \cite{Uhlenbeck1982Removable} with the bounded inhomogeneous
terms retained.

Compare the Chern integrals of the original sequence and this extension
on a surrounding sphere collar of strong $W^{1,p}$ convergence, with
$p>4$.  Fubini's theorem selects a boundary sphere on which a subsequence
converges in tangential $W^{1,p}$; the Chern--Simons boundary integrals
then converge.  More precisely the difference of the two ball charge
integrals is the degree of their determinant-one clutching map
$g_j:S^3\to\SU(2)$ plus the difference of the boundary
Chern--Simons integrals in the convergent gauges.  The latter tends
to zero.  Bounded charge bounds the integer degrees, so a subsequence
has constant degree and the limiting charge difference is that integer.
Since the self-dual part has no
atom, Chern--Weil identifies the lost curvature mass with $8\pi^2$
times that integer.  The mass is positive, so the integer is positive.
Previously labelled weights simply add to new weights; collisions give
the same cutoff expression because the summands $m_j\chi(x_j)$ add.
The limiting equations on every fixed position stratum are therefore
smooth, particle-dependent equations of the stipulated form.

All arguments are uniform for coefficients converging in the weighted
seminorms: approximate by a finite library, use the finite convergence
argument, and use the uniform tail bounds.  If no further degeneration
occurs, the full bootstrap and the true-end estimate below give
convergence in the completed field charts.  This proves the lemma.
\end{proof}

\begin{lemma}[True necks and exponential completion]
\label{lem:analysis:neck}
Along a sequence in which genuine neck lengths tend to infinity, a
subsequence has finitely many particles and nonconstant cylinder levels.
Every translated cylinder level has zero spinor and is ASD.  Nonconstant
levels have a discrete positive charge cost.  The remaining necks carry
no diffuse loss and converge to matching flat states with exponential
decay of all derivatives.  On a no-loss stretch of length $T$ the
deviation from the flat orbit is bounded by exponentially decreasing
contributions from its two ends and from the prescribed forcing.
\end{lemma}
\begin{proof}
Deep translation eliminates the background splice errors and all
nonstationary graph terms.  The limit is an unperturbed solution on a
positive-scalar-curvature cylinder with flat fixed-line background.
The integrated Weitzenb\"ock inequality with cutoffs and finite energy
forces its spinor to be zero.  It is therefore ASD.  We first prove
the small-strip estimate and completion of any one such finite-energy
level, without assuming a finite extraction count.

The flat orbits on $S^3$ and on the lens space form a finite list.
For a lens this follows by diagonalizing the representation of its
finite cyclic fundamental group.  Their adjoint $H^1$ vanishes;
one may also obtain this by averaging over the finite covering group.
Positive scalar curvature kills the flat Dirac kernel.  Thus their
normal boundary operators have a common positive gap $\beta$.
On any stretch close to a flat $\gamma$, put the spatial connection
$a=\gamma+\xi$ in the slice $d_\gamma^*\xi=0$.  Use the parallel
group to remove the parallel component of the temporal connection.
The remaining temporal component $\alpha$ is determined, on the
complement of constants, by an equation of the form
\[
 d_\gamma^*d_a\alpha=d_\gamma^*\bigl(Q(\xi,\Phi)+f_+\bigr),
\]
where $Q$ is quadratic.  Its inverse from $H^{-1}$ to $H^1$ gives
$\|d_a\alpha\|_2\leq C(\|Q\|_2+\|f_+\|_2)$, uniformly for
small $\xi$.  Hence the projected nonlinear term is of order zero
in the low norm.  Writing $u=(\xi,\Phi)$, the projected equation is
\begin{equation}\label{eq:analysis:neck-evolution}
 (\partial_t+B_\gamma)u=N(u)+f(t),
 \qquad \|N(u(t))\|_{L^2(Y)}
             \leq C\varepsilon\|u(t)\|_{L^2(Y)}.
\end{equation}
The transformed forcing $f$ has the same strip Sobolev bounds as the
original forcing.  The operator $B_\gamma$ is the co-closed
$*_3d_\gamma$ block together with the flat Dirac block; all parallel
gauge constants have been separated.  Higher-order tame estimates
give $\|N(u)\|_{H^{k-1}(I)}\leq C_k\varepsilon\|u\|_{H^k(I)}$
on a fixed-size strip, with a fixed enlargement of $I$ if required.
Smallness in the needed local norms follows from small-energy
regularity on these genuine collars.

Here are the norms in the decay estimate.  Move the ends by less
than one, absorbing these bounded collars into the end norms, so
the length of the stretch is an integer $M\geq6$.  Put
$I_j=[j,j+1]\times Y$ and
$I_j^+=[\max(0,j-2),\min(M,j+3)]\times Y$, for $0\leq j<M$.
Set
\[
 A_j=\|u\|_{H^k(I_j)},\qquad
 F_j=\|f\|_{H^{k-1}(I_j^+)},\qquad
 A_-=\|u\|_{H^k([0,3]\times Y)},\quad
 A_+=\|u\|_{H^k([M-3,M]\times Y)}.
\]
For any sufficiently small $0<\delta<\beta$ one has
\begin{align}\label{eq:analysis:neck-decay}
 A_j\leq{}&C_k\left(e^{-\delta j}A_-
                 +e^{-\delta(M-1-j)}A_+\right)\notag\\
          &+C_k\sum_{r=0}^{M-1}e^{-\delta|j-r|}F_r.
\end{align}
Its constants are independent of $M$.  To prove it, first propagate
the positive spectral subspace of $B_\gamma$ from the left and the
negative subspace from the right in $L^2(Y)$.  The two semigroups
have $L^2$ operator norm at most $e^{-\beta|t-s|}$.  Integrating over
unit strips and using Cauchy--Schwarz yields, with
$L_j=\|u\|_{L^2(I_j)}$, the discrete bound
\[
 L_j\leq C\left(e^{-\beta j}A_-+e^{-\beta(M-1-j)}A_+\right)
       +C\sum_r e^{-\beta|j-r|}(\varepsilon L_r+F_r).
\]
End traces may be chosen in the first and last unit strips by
Fubini, or bounded by $A_\pm$.  The exponentially weighted sequence
norm with rate $\delta<\beta$ bounds this convolution uniformly in
$M$, so sufficiently small $\varepsilon$ absorbs its $L_r$ term.

Next, the interior elliptic estimate for $\partial_t+B_\gamma$ on
overlapping strips gives
\[
 A_j\leq C_k\left(\|u\|_{L^2(I_j^+)}+F_j
                   +\varepsilon\|u\|_{H^k(I_j^+)}\right)
\]
away from the two end collars.  Only a fixed number of neighboring
$A_r$ enter the last norm.  In the same exponentially weighted
sequence norm this finite-neighbor operator is bounded uniformly in
$M$; absorb it after further decreasing $\varepsilon$.  Substituting
the low-norm bound proves \eqref{eq:analysis:neck-decay}, with a
possibly smaller $\delta$.  This two-stage argument uses a nonsingular
$L^2$ semigroup estimate and then elliptic regularity.  It does not
assert a bounded $H^{k-1}\to H^k$ semigroup kernel at $t=s$.
Increasing $k$ gives all required spatial and temporal derivatives.

Now consider either tail of one finite-energy ASD cylinder level.
Its unit-strip energies tend to zero.  Local small-energy compactness
therefore puts all sufficiently remote strips near the finite set of
flat orbits.  Overlapping strips and continuity force an entire tail
to remain near one orbit, since distinct sufficiently small flat
neighborhoods are disjoint.  The slice and estimate just proved apply
on arbitrary finite subintervals of that tail.  Sending their remote
endpoint to infinity gives exponential convergence to the flat;
boundedness of its end norm follows from small-energy regularity.
Only now apply the Chern--Simons identity to its completed ends.
The finite flat list has discrete Chern--Simons differences modulo
integer charge, and a nonconstant ASD level has positive energy.
Consequently these levels have a uniformly positive charge cost.
Together with the integral particle masses and
\eqref{eq:analysis:global-energy}, this bounds the extraction count.

After those finitely many extractions every remaining unit strip is
close to a flat orbit; a violating strip translated to the origin
would produce a further level or particle.  The same slice argument
and \eqref{eq:analysis:neck-decay} apply to each residual stretch.
They keep it in one flat neighborhood and give exponential decay
from its two ends, plus the exponentially weighted forcing sum.

There is no diffuse charge in this residual stretch.  In the slice
neighborhood choose one lift of Chern--Simons; its difference between
two slices tending to the same flat tends to zero.  This is the
topological charge between them.  The positive local Weitzenb\"ock
estimate, or \eqref{eq:analysis:neck-decay}, controls the spinor and
self-dual energy there.  Chern--Weil then controls the remaining
curvature energy.  Thus no energy is missed by the discrete extractions.
Exponential forcing gives exponential completion and, in a no-loss
middle, exponential smallness.  Comparisons by the compact residual
parallel groups at the ends are allowed by our gauge convention.
The forcing bounds are uniform under arbitrary simultaneous approaches
to true faces by \eqref{eq:analysis:exponential-seminorm} and the
product collars.  The removed parallel temporal component is a pure
gauge mode, so a slowly varying comparison of end constants cannot
produce a residual polynomial forcing term.

This lemma concerns isolated rational flat states.  The harmonic
flat-slice argument at an ordinary $S^1\times S^2$ auxiliary neck is
the separate one in Section~\ref{sec:ordinary}; no assertion about
that positive-dimensional flat variety is being deduced from the
spectral gap used here.
\end{proof}

\begin{lemma}[Exact support and incidence budget]
\label{lem:analysis:incidence}
At an ideal limit a basic sphere, surface, or point test persists unless
a particle meets its exact basic support.  In that case one may drop
the test with its auxiliary evaluation variable, while retaining the
particle's incidence with that support.  For every particle of weight
$m$, the dimension of its permitted position plus the degrees released
by all such dropped tests is at most $4m$.  The zero-evaluation sphere
rules in Definition~\ref{def:indices:analytic-hypotheses} hold, with
different charged tests using distinct loss units.
\end{lemma}
\begin{proof}
If no particle meets the basic sphere support, the sphere path converges in
$C^{0,\beta}$ with $\beta>b$, hence in the $H^b$ topology needed in
Lemma~\ref{lem:analysis:zero-winding}.  The modified sweep and every
surviving correction converge.  A correction graph hitting an atom
vanishes by \eqref{eq:analysis:energy-cutoff}; this does not invalidate
the limiting section formula.  Thus the test persists unless its basic
sweep support itself is hit.  At a $J$ face the sweep has already been
segmented; the mark lies on exactly one nonidentity segment and assigns
the event to that side.  Identity segments are uniformly disjoint from
the event and cannot be limits of its marks.  Cap sections at an
unhit evaluation point have the same convergence and cutoff argument.
Phase sections and all other sampled cut terms converge by their
cutoff formulas; at a fixed tuple phase cuts are not needed for the
fixed-type exclusion.

The position and released-degree counts are explicit.  A sphere or
surface hit gives at most two position parameters and releases degree
two.  For a possibly degenerate sphere map, retain a preimage in its
compact two-dimensional parameter domain; injectivity is unnecessary
for this upper bound.  A fixed point gives no position parameter and
releases four degrees.  Two distinct cap surfaces meet transversely,
so a simultaneous hit has zero position dimension and releases four
degrees.  There are no triple incidences, competing sphere supports,
or sphere--cap incidences.  An unconstrained particle has four position
parameters and releases no degree.  These cases give at most four for
weight one; multiplicities only increase the allowed bound $4m$.
The lower-rank matrix stratum has higher codimension and cannot increase
this budget.

For a parallel field, a whole test with zero evaluation is impossible
without a support hit by Lemma~\ref{lem:analysis:zero-winding} and the
absolute smallness of its variations.  During segmentation the same
conclusion applies when both half evaluations vanish, by
Lemma~\ref{lem:analysis:segmentation}.  At the split it applies to the
assigned half separately.  At a lens face the whole test lies in $N$;
the charge-zero extendibly central reference has the same avoidance,
and the index and charge table of
Proposition~\ref{prop:indices:lens-index} charges any remaining permitted
alternative to the stated excess.  Disjoint basic sphere supports
prevent one particle unit from paying for two such tests.  Assign a
unit to each witnessing particle on a main support, and assign the
necessary excess in a tested lens cap to that cap's single sphere.
The different caps are disjoint.  Near a $J$ separation each incident
particle remains on the intrinsic compact support of its assigned
half, so it cannot escape into the connector while carrying a second
test.  A particle carried into a newly separated $N$ transfers its
integer weight to that cap's excess; when its evaluation is zero the
reference charge is zero, so this weight still pays for its one test.
Unconstrained particles escaping to a connector carry their charge
there but create no additional zero-degree support incidence.
Colliding labels add their weights and never identify two charge units.
These assignments, with each connector charge allocated once in the
virtual decomposition, prove the distinct-unit assertion.

In particular a main monopole particle removes six field dimensions
and restores at most four through incidence and position variables.
Its net penalty is at least two.  The ordinary instanton loss is eight,
so its corresponding penalty is at least four.
\end{proof}

\subsection{Candidate main fields and their degenerations}
\label{subsec:analysis:candidates}

The dimension argument of Section~\ref{sec:indices} will be applied to
a larger class than the actual ideal limits.  We keep the equations on
every main component but replace the lens-cap and cylinder fields by their
possible charge assignments and impose no matching across true cuts.
The retained equations will give ideal compactness for these main fields;
the assigned losses will bound the number of degenerations.  Neither
assertion will require reconstruction of the omitted fields.

A \emph{candidate} consists of completed main fields, their parameter
values, weighted main particles, retained cut variables, and a discrete
charge assignment.  The main fields solve their equations and all
retained nonphase cuts for the current perturbation data.  When phase
sections have been installed, we also impose those cuts.  For
single-valued data this is a single system of equations; after finite
branch lists have been added, a candidate solves one local lift of that
library.  The class includes every actual ideal boundary tuple for the
same data, but its omitted lens caps and matchings need not be reconstructible.

The charge and incidence data are as follows.  At a $J$ connector
record its additional nonnegative integer charge; at a lens record an
effective cap charge in the range and congruence determined by the
main flat state and its required sphere test.  Retain the total charge
condition.  At a labelled point on a basic support allow the
corresponding test and its auxiliary evaluation variable to be dropped,
but retain its exact support incidence
as in Lemma~\ref{lem:analysis:incidence}.  There are finitely many
assignments of a segmented test to a side at each $J$ face.
At a finite switching value use the closed support charts and limiting
union specified after Lemma~\ref{lem:analysis:segmentation}; near a true
$J$ face only the intrinsic half-support charts occur.  Include the
boundary strata of these closed parameter ranges among the finite
support-stratum choices.

We first bound the candidate fields and charges.  The main charges
need not be nonnegative, so the total charge condition alone is not
an energy bound for each piece.  Apply the integrated Weitzenb\"ock
inequality on each completed main piece first, using
Young's inequality on its negative zeroth-order terms.  Bounded
squared costs of those terms and of the forcing give uniform bounds
on $\|\Phi_i\|_4^4$ and $\|\nabla\Phi_i\|_2^2$, independently of
the charge.  The curvature equation then gives
$\|F_i^{0,+}\|_2^2\leq B_i$.  Write
$\kappa_i^{\rm main}$ for the actual charge of this smooth completed
main field, excluding particles and omitted end fields; this is
distinct from the virtually filled $\kappa_i$ of
Section~\ref{sec:indices}.  Thus
\begin{align}\label{eq:analysis:candidate-energy-order}
 \kappa_i^{\rm main}&\geq-\frac{B_i}{4\pi^2},\qquad
 \mathcal L\leq\kappa^{\rm tot}+\sum_i\frac{B_i}{4\pi^2},\notag\\
 \kappa_i^{\rm main}&\leq\kappa^{\rm tot}+
                         \sum_{j\ne i}\frac{B_j}{4\pi^2}.
\end{align}
Here $\kappa^{\rm tot}=\sum_i\kappa_i^{\rm main}+\mathcal L$ and
$\mathcal L\geq0$ is the omitted effective end and particle loss.
Finally
$\|F_i^0\|_2^2\leq2B_i+8\pi^2\kappa_i^{\rm main}$ gives the full curvature
bound.  The finitely many component and face types have uniform
constants $B_i$.  This argument requires only the completed main
equations and the assigned nonnegative losses.  It supplies the energy
bound for the degeneration argument that follows.

\begin{lemma}[Finiteness of degeneration data]
\label{lem:analysis:finite-data}
For this finite problem the degeneration data needed for candidate
exclusion have finite ranges.
Order the data by: deeper true faces;
larger total effective end and particle loss; fewer distinct main
particle locations at fixed loss; more dropped tests; and more $R$
components.  Treat finite assignment, closed support-stratum, and rank
choices together at each step.  In a fixed finite branch library, if the single-valued
coefficients converge in $\mathfrak P$ and the additional finite
coefficients converge, every sequence of candidates has a subsequence
which either converges modulo the completed field slices at the same
step or limits to an earlier step or a fixed-type tuple.
\end{lemma}
\begin{proof}
There are finitely many true faces and flat-state assignments.  The
energy bound gives a uniform upper bound on all nonnegative end
losses and particle weights.  These charges lie in discrete sets
(quarters suffice for the lens data and integers for particles and
$J$ losses), so their ranges and the possible numbers of positions
are finite.  There are finitely many basic tests and components.

Relative bundle sectors are also finite after inessential boundary
degree relabellings.  With determinant fixed and limiting flats fixed,
relative $c_2$ is an integer and is bounded by Chern--Weil and the
candidate curvature bound just proved.  For split fields, every possible
completed core has rational-homology-sphere boundary, so compactly
supported two-forms detect all of its real second cohomology.  A finite
dual collection $\beta_j$ satisfies
\[
 |\langle v,[\beta_j]\rangle|
 \leq (2\pi)^{-1}\|F_L-F_M\|_2\|\beta_j\|_2.
\]
Thus the full curvature bound bounds every coordinate of the integral
splitting lattice; torsion gives only finitely many further choices.
The finitely many split cores and compactly supported representatives
are fixed in the current metric problem.  Holonomy gauges of nonzero
degree only relabel the relative sectors with the same stated charge
data.  These observations give the asserted finite ranges.

Apply body compactness and neck extraction.  At a newly appearing cut,
Lemmas~\ref{lem:analysis:compactness} and \ref{lem:analysis:neck}
give the stated charge assignments.  True cuts persist, and a
new cut is an earlier step.  A new bubble or positive cylinder level
increases the total loss.  An already separated lens may acquire a
different main asymptote only through an extracted cylinder level.
If the old state is $\gamma$, the new state is $\gamma'$, and the
extracted ASD cylinder runs from $\gamma'$ to $\gamma$ with charge
$e>0$, replace its assigned cap charge by
\begin{equation}\label{eq:analysis:cap-charge-transfer}
 \kappa_N'=\kappa_N+e.
\end{equation}
The Chern--Simons identity gives exactly the charge congruence for
$\gamma'$, with the fixed determinant transitions.  The previous
tested cap charge is positive.  Hence the new charge is at least
$1/4$ if the new state is trace, and at least one if it is central,
by its respective congruence.  It satisfies the same allowed range
and lies at an earlier loss step.  This argument needs the cylinder
extracted from the main solution, not a reconstructed old cap.
Several levels add their charges; a pure parallel-frame change has
zero charge and does not change the flat state.

Escape of an already labelled particle into
an already separated end transfers its weight to effective end loss;
if the total loss stays constant, the number of main locations
decreases.  Collisions also decrease that number.  A surviving basic
test can fail only by an exact hit, which adds a dropped test; existing
hits persist by closedness of their incidence parameterizations on
that face.  In particular an incident particle cannot escape down a
true $J$ middle: the rule has already switched there, and the assigned
half test has a compact intrinsic support in its retained component.
It remains incident on that half in the limit.  At a newly appearing
lens face the whole test lies in the separated $N$ core and is
omitted from the main problem altogether; a particle incident to it
transfers its integer weight to that cap charge.  A decrease in the
number of dropped \emph{main} tests here is harmless, since the
new true face is earlier in the ordering.  At an already separated
end every retained main test has compact support, so only particles
without such an incidence can escape down that end.  Thus no missing
main test is restored without justification.  Earlier whiskers used
during a switching homotopy are absent in the neighborhood of the
true face and cannot carry an incidence for its current test.
The marks cannot escape through identity segments of sphere sweeps.
The $R$ condition is closed under smooth completed convergence: a
normalized parallel splitting has a convergent subsequence in the
compact space of base lines and remains parallel, and the spinor
remains zero.  A new $R$ component therefore moves earlier.

If none of these worsenings occurs, there is no new body or end loss.
The strong bootstrap and exponential completion give convergence in
the same high-order field charts.  Changes among $A,S,P$ cause no
singularity in the full non-$R$ field model as long as at least
one $P$ remains: the finite-isotropy construction and compatible
charts of Lemmas~\ref{lem:analysis:samples} and \ref{lem:analysis:scale}
apply throughout.  If the last $P$ disappears, the limit is a fixed-type
tuple.  Rank-one tests may limit to rank zero; cover them using full
matrix-value neighborhoods and impose transversality to each rank
stratum there.  This avoids an unjustified compactness assertion for
the open rank-one stratum alone.
\end{proof}

\subsection{The fixed stage and legitimate instanton links}

\begin{proposition}[First-stage regularity]
\label{prop:analysis:first-stage}
There is an arbitrarily small single-valued, single-component choice
for which type $A$ equations with all retained nonphase tests are
transverse on every relevant ideal and parameter stratum, and type
$S$ tangential equations are transverse.  The first two conditions of
Definition~\ref{def:indices:analytic-hypotheses} follow.  The only
fixed-type tuples are a finite set of legitimate unbroken cut-down
instantons, and the coupled complex Dirac operator can also be made
surjective at each of them without altering the instanton equations.
\end{proposition}
\begin{proof}
Use the countable single-valued library with the locality and bit
identifications imposed above.  Lemma~\ref{lem:analysis:samples} and
the universal-surjectivity argument give regularity on each $A$
equation and retained-cut problem and on each pure $S$ tangential
problem.  There are countably many relative topologies, coordinate
charts, and finite position and rank strata.  The intersection of
their residual regularity sets in the arbitrarily small allowed open
ball of $\mathfrak P$ is residual.  At a true $J$ face, a component's
equations use only its retained parameters, so this assertion does not
give it spurious parameter directions from another component.

For $A$, the index after the reference fillings is $q_i$.  Each assigned
effective lens charge supplies the correction $\delta_I$, while removal
of its one parameter
and degree-two test restores one, giving the penalty $\delta_I-1$.
The permitted cap charge and test exclude the zero-charge central
case, and all remaining lens penalties are positive by
Proposition~\ref{prop:indices:lens-index}.  A main instanton particle
has net penalty at least four by Lemma~\ref{lem:analysis:incidence},
and an allocated positive integer $J$-trajectory charge has positive
index penalty.  A regular nonempty problem therefore has $q_i\geq0$,
strictly if any of these losses occurs.

For $S$ retain only the pure tangential equations.  The active line is
ordered by its nonzero spinor.  Use the sharp reference fillings of
Proposition~\ref{prop:indices:lens-index}, whose tangential correction,
including the limiting line stabilizer, is zero.  The actual pure-line
problem has index $d_s+p_i$, with the prescribed ideal labels adding
at most $4\ell_i$ position directions when tests are ignored.
Transversality implies $d_s+p_i+4\ell_i\geq0$.  This is precisely the
tangential bound required; it makes no assertion of full normal
regularity at an arbitrary $S$ field.  The incidence condition is
Lemma~\ref{lem:analysis:incidence}.

Apply Proposition~\ref{prop:indices:fixed-exclusion}, which uses only
these fixed-stage bounds and the metric clamps.  Every fixed tuple
except an unbroken, loss-free, cut-down type $A$ field is excluded.
Compactness from Lemmas~\ref{lem:analysis:compactness} and
\ref{lem:analysis:neck} now makes that zero-dimensional regular
instanton set compact, hence finite.  This reasoning uses compactness
of sequences and the already proved fixed-type bounds; it does not
assume compactness of the final monopole bordism.

At such an instanton $a$, phase equivariance forces $P_D(a,0)=0$,
but its spinor derivative may be any complex-linear compact operator.
Let $K_a$ and $C_a$ be the kernel and cokernel of its current complex
Dirac operator.  Finitely many transported spinor samples inject
$K_a$, and finitely many localized smooth output spinors detect
$C_a$.  Hence their compositions give every map $K_a\to C_a$.
The irreducible connection has only its central gauge stabilizer;
complex-linear spinor operators respect both that center and phase.
They extend from its sample slice with compact cutoffs.  Since
$\dim_{\C}K_a-\dim_{\C}C_a=n_D^{\rm tot}>0$, a generic arbitrarily
small such variation makes the Dirac operator surjective.  The value
at zero spinor remains zero and the ASD equation and its cuts are
unchanged.  Treat the finitely many instantons in disjoint configuration
neighborhoods.

These additional finite conditions are open.  They can be combined
with the preceding residual choices; alternatively make the small
finite changes and then choose generically within their open
surjectivity neighborhoods.  Absence of the other fixed tuples persists
under sufficiently small changes by ideal compactness.  Existing true
cuts and positive losses persist in a limit.  An unbroken $R$ or $S$
sequence retains a parallel reduction and cannot tend to a legitimate
irreducible instanton.  The finitely many legitimate instantons have
regular isolated ASD neighborhoods, and every nearby instanton lies
in their union, since an outside sequence would have an excluded
ideal limit.  Their Dirac surjectivity is an open condition there.
Consequently the finite adjustments and the subsequent generic choice
preserve the fixed-type exclusion.  This completes the first stage.
\end{proof}

Fix this ultimate first-stage choice before constructing the quadratic
phase sections below.  Its finitely many instantons and their
surjective Dirac kernels are therefore the centers of those charts.
Subsequent mixed and interior regularizing variations are identically
zero in smaller fixed neighborhoods of the instanton collars.  They
do not move the centers or their quadratic leading sections.  Thus
no assertion that genericity on every ideal component stratum is open
is needed: after each preliminary adjustment genericity is reselected
inside the open neighborhood preserving the forbidden-tuple absence,
the finite ordinary zeros, and their Dirac surjectivity; the link data
are installed only after that selection is complete.

\begin{proposition}[The instanton link and its phase degree]
\label{prop:analysis:instanton-link}
Put $n_*=n_D^{\rm tot}>0$ and $\eta=n_*-1$.  At each legitimate
cut-down instanton, the normal spinor link before phase quotient is
$S^{2n_*-1}/\{\pm1\}$ and its phase quotient is $\CP^{n_*-1}$.
The phase-character-two line restricts to $\mathcal O(2)$ on this
projective space.  The $\eta$ phase sections can be chosen so that
their cut-down link is regular and has signed count
\begin{equation}\label{eq:analysis:instanton-degree}
 2^\eta\,\operatorname{sign}(a).
\end{equation}
There is one fixed overall boundary-orientation convention in this
formula, the same for every instanton.  Subsequent mixed branch terms
may be kept zero near these links.
\end{proposition}
\begin{proof}
The cut-down ASD and metric problem is a regular isolated point and
the complex Dirac operator is onto.  Its kernel $K$ has dimension
$n_*$.  The implicit theorem, with the cut marks included among the
finite variables, solves the transverse field, metric, and cut
coordinates in terms of the small spinor kernel.  Curvature outputs
are phase-invariant, so their first spinor derivative at zero is zero;
Dirac outputs have phase weight one.  The normal model is consequently
a smooth equivariant deformation of $K$, with the gauge center
identifying $\psi$ and $-\psi$.  Equivalently write
$\Phi=r\psi$, $\|\psi\|=1$, and divide the Dirac equation by $r$.
At $r=0$ it is the surjective kernel problem, and radial blow-up gives
boundary $S(K)/\{\pm1\}$.  Dividing by phase gives $\mathbb P(K)$.
Its orientation is the complex orientation multiplied by the ordinary
instanton sign.

Choose phase sections quadratic to leading order in $\Phi$.  Finitely
many spinor samples inject the finite-dimensional space $K$, so any
complex quadratic form on $K$ extends to a quadratic function of these
samples.  The connection sample slice makes this extension equivariant;
the center acts trivially on squares.  Put the support in an interior
parameter chart near the instanton.  Division by $r^2$ extends its
section smoothly over the blow-up and gives the chosen quadratic form
on the projective boundary.  The freedom is independent of the other
marks and of the equation perturbations.

More explicitly, a homogeneous function $q(z)$ of phase weight two
satisfies $q(\lambda z)=\lambda^2q(z)$ and defines a section of
$\mathcal O(2)$ on $\mathbb P(K)$.  The center quotient does not turn
this into $\mathcal O(1)$: the principal bundle
$S(K)/\{\pm1\}\to\mathbb P(K)$ has doubled Hopf transition
functions when its effective circle is identified with a standard
circle.  Thus its associated weight-two line for the original phase
has $c_1=2h$.  This agrees with the link calculation in
\cite[Lemma~3.28 and Proposition~3.29]{FeehanLenessPU2II}.

For a direct regular choice, in coordinates
$[z_0:\cdots:z_{n_*-1}]$ take
$z_j^2-a_jz_0^2$, $1\leq j\leq n_*-1$, with all $a_j\ne0$.
Every common zero has $z_0\ne0$, there are exactly $2^{n_*-1}$
such points, and their complex Jacobian is invertible.  This also
proves the intersection number
$\langle(2h)^{n_*-1},[\CP^{n_*-1}]\rangle=2^{n_*-1}$.
For $n_*=1$ there are no phase cuts and the same formula is one.
The implicit theorem extends these boundary points to regular collars.
Choose every later coupled term outside smaller neighborhoods of these
finitely many collars.  Where such a term is zero its repeated branches
have total weight one, so no extra rational factor changes
\eqref{eq:analysis:instanton-degree}.
\end{proof}

\subsection{Mixed projections and finite exclusion}

\begin{lemma}[The necessary mixed projection]
\label{lem:analysis:projection}
For a candidate containing a $P$ field, delete the $R$ fields, their
particle positions, and their tests, but retain their parameter
variables.  Quotient the remaining fields by their independent gauges
and the common phase.  This defines a Fredholm problem with finite
isotropy near its candidate projections.  Its dimension, before phase
quotient and phase cuts, is at most
\begin{equation}\label{eq:analysis:projected-dimension}
 \sum_{i\notin R}\left(
 i_i-\sum_{\text{lens ends of }i}(\delta_{\rm sp}-1)-2w_i\right)
       +\sum_{i\in R}p_i.
\end{equation}
The phase cuts and quotient subtract $2\eta+1$.  At every such
candidate projection, admissible local finite-branch variations make
the universal equations and retained cuts transverse.
\end{lemma}
\begin{proof}
Lemma~\ref{lem:analysis:R-independence} proves that the equations
factor through this projection, including every energy-cutoff atom.
The $P$ witness makes its isotropy finite by
Lemma~\ref{lem:analysis:samples}.  Particle positions satisfying a
dropped sphere incidence can be parameterized by preimages in the
sphere domain; a surface incidence uses its surface domain, and two
surface incidences use their transverse fiber product.  These are the
finite-dimensional position spaces used in computing the index.
Surviving test marks are retained separately.  Singularities of a
possibly nonimmersed sphere map do not increase that parameter-space
dimension.

The virtually filled field index is $i_i$.  Removal of a lens filling
costs $\delta_{\rm sp}$, while omitting its one parameter and degree-two
test restores one.  Its penalty is therefore the full
$\delta_{\rm sp}-1$, not merely one per lens.  Every unit of main
particle charge loses six field dimensions and recovers at most four
through the position and incidence budget.  Allocated integer
connector losses lose at least as much as the bound two charged here.
The discarded $R$ equations contribute no fields or cuts to the
projection, only their $p_i$ remaining parameters.  This proves
\eqref{eq:analysis:projected-dimension}; higher-codimension rank strata
only decrease it.  Finally the common phase has finite stabilizer and
the $\eta$ complex phase sections impose $2\eta$ real conditions.

Full equation outputs at every retained target, full matrices at cap
tests, holonomy-event values, and independent phase-section values
are available by Lemma~\ref{lem:analysis:samples}.  The
universal-surjectivity argument and Lemma~\ref{lem:analysis:scale}
apply on each finite branch lift.  No transversality statement on an
arbitrary removed $R$ field is used.
\end{proof}

\begin{proposition}[Finite exclusion of mixed candidates]
\label{prop:analysis:finite-induction}
After the first-stage choices, finitely many additional local
multisection terms with arbitrarily small coefficients exclude all
forbidden mixed candidates.  The only remaining mixed candidate types are
the unbroken free part and the single trace-minimum lens-end main
problem of Proposition~\ref{prop:indices:mixed-exclusion}.  The latter
is regular of dimension zero.  The unbroken free problem has expected
dimension one; its regularization follows the lens gluing in
Section~\ref{sec:gluing}.  The additional choices agree on paired lens
exteriors and are zero near the legitimate instanton collars.
\end{proposition}
\begin{proof}
The inequalities in the proof of
Proposition~\ref{prop:indices:mixed-exclusion}, applied to
\eqref{eq:analysis:projected-dimension}, give strictly negative
projected index for each forbidden candidate type.  Here the numerical
argument is used to compute the index, before asserting that a
negative index eliminates a zero.  The preceding fixed stage already
excluded all unwanted fixed-type tuples.  We now establish the
regularity that turns these negative indices into absence.

Use the finite ordering of Lemma~\ref{lem:analysis:finite-data}.
Suppose all earlier steps have been excluded at the current
coefficients.  The set needed at this step is the image under
projection of the candidates of Section~\ref{subsec:analysis:candidates},
with their $R$ fields still present before projection.
Lemma~\ref{lem:analysis:finite-data} makes this image compact modulo
local slices and its finite assignment choices: a sequence has a limit
at this step, at an already excluded earlier step, or at an already
excluded fixed tuple.  It is not necessary, and is not asserted, that
every zero of the loosened equations after deleting the $R$ fields
belongs to a compact set.

A forbidden candidate cannot instead converge to a legitimate clean
instanton.  Any existing cut or positive loss persists in the
degeneration order.  It cannot become the single permitted trace end
either: already separated end loss cannot decrease, and a particle
escaping into a newly separated lens carries at least one unit,
whereas the permitted total lens charge is $1/4$.  The unbroken free
stratum with no degeneration is not one of the exclusion steps.
These facts justify the use of the already excluded fixed locus and
of protected neighborhoods of the legitimate instanton collars.

At each point of the compact candidate image choose finitely many
supported variations from Lemma~\ref{lem:analysis:projection}
spanning the missing universal directions.  Take a finite subcover
and the finite sum $V$ of those coefficient spaces.  There are only
finitely many previously chosen branch lists.  At each local lift
include all their branch choices and all translates in the newly
chosen coefficient spaces.  The invariant-span prescription in
Lemma~\ref{lem:analysis:samples} makes the universal map transverse
for every branch separately.  At zero new coefficients a new list
merely repeats the old list with total weight one.

Surjectivity persists near this compact image for sufficiently small
$v\in V$.  Moreover every new candidate for sufficiently small $v$
has its projection in this neighborhood.  Otherwise take $v_j\to0$
and candidates outside it.  Compactness gives either a forbidden
earlier limit or a candidate at the current step whose projection is
inside the neighborhood, both contradictions.  The same argument
shows that absence at all completed steps is open under sufficiently
small changes.

Apply finite-dimensional parametric transversality to $V$, in the
countably many projected field charts and to each local branch and
rank stratum.  The universal zero set is smooth by
Lemma~\ref{lem:analysis:scale}; the regular parameter values are
residual (equivalently apply Sard--Smale to its parameter projection).
At any such value, every projected negative-index zero set in the
chosen neighborhood is empty.  Since every candidate projects there,
the current step is excluded.  Choose the coefficients small enough
to preserve all previous exclusions, the metric tolerances, and the
protected instanton collars.  There are finitely many steps, each
adding finitely many terms, so the total product of branch choices is
finite.  This turns every negative projected index used in
Proposition~\ref{prop:indices:mixed-exclusion} into absence, proving
its exclusions.  Regularity of the remaining strata is a separate step.

For the single trace-end main problem, use the same construction to
obtain regularity instead of emptiness.
Its compact candidate set is isolated from forbidden corners by what
was just proved.  It has a single type $P$ main component; its effective
sample action permits ordinary single-valued variations.  Choose these
simultaneously on the paired lens exteriors, using their exact tensor
identification.  Its expected dimension is zero.

All additions use the corner and bit-flip extension prescriptions,
so the required exterior equalities are preserved.  This completes
the forbidden-candidate exclusion and trace-main regularization.
\end{proof}

\begin{lemma}[The weighted boundary rule]\label{lem:analysis:stokes}
A compact oriented one-dimensional zero set built from the finite
branch lists above in the surviving free stratum has total weighted
oriented boundary zero.  Its
weights are rational, and a region on which every additional branch
term is zero has the weight of the original single-valued problem.
\end{lemma}
\begin{proof}
In this application the unbroken stratum has one $P$ component.
After quotienting the ineffective diagonal center, its gauge and
phase action is free.  Thus its local base charts are ordinary
manifold charts; no additional effective-orbifold integration factor
is being omitted.
In a sample slice tube every $H$ translate has weight $1/|H|$.
Changing the slice lift permutes these translates and preserves their
multiplicities.  To check orientations, choose the finite samples to
detect the exact combined stabilizer at the central candidate.  The
$P$ witness restricts phase to $\{\pm1\}$.  A nonzero-spinor target
with such a phase has scalar parallel gauge $\pm1$: in $\SU(2)$
an eigenvalue $\pm1$ forces both eigenvalues to be that scalar.
An $A$ target already has only central gauges.  Thus $H$ is a
subgroup of a product of central groups $\{\pm1\}$, with the
common phase, modulo its ineffective kernel.  Nearby tube stabilizers
are conjugate to subgroups of $H$ and are still central.  They act
trivially on the real adjoint, parameter, and cut directions and
complex-linearly by signs on the spinor blocks.  Deform the
linearization equivariantly to its real ASD part and complex Dirac
part.  These actions preserve the determinant orientation there,
and hence preserve it throughout the deformation.  Each cutoff is
zero on a neighborhood of the tube boundary, so extension by repeated
zeros creates no boundary of a branch zero set.  Finite products have
product weights, and the sum of their weights is one.  Consequently
the weighted oriented germs agree on chart overlaps, even where
branches coincide or separate.

Choose a partition of unity on a finite cover of a compact remainder
and apply the one-dimensional Stokes formula on each oriented local
lift with its weight.  Terms involving derivatives of the partition
sum to zero because the weighted germs agree.  The remaining terms
are exactly the weighted boundary points.  This proves their total
is zero.  At a zero perturbation all repeated branches represent the
same germ and their total weight is one, proving the final assertion.
\end{proof}
